% Lecture notes template for the /lecture skill (fixed format).
% Replace every <...>. Keep only the sections this topic needs; copy macros from previous lectures.
\documentclass[11pt,a4paper]{article}
\usepackage[margin=2.3cm]{geometry}
\usepackage{amsmath,amssymb,amsthm,bm}
\usepackage{booktabs}
\usepackage{enumitem}
\usepackage{xcolor}
% Only if the lecture has an algorithm:
% \usepackage{algorithm}
% \usepackage[noend]{algpseudocode}
\usepackage{hyperref}
\hypersetup{colorlinks,linkcolor=blue!50!black,urlcolor=blue!50!black}

\newtheorem{theorem}{Theorem}
\newtheorem{lemma}[theorem]{Lemma}
\theoremstyle{definition}
\newtheorem{definition}{Definition}
\newtheorem{example}{Example}
\theoremstyle{remark}
\newtheorem*{remark}{Remark}

\newcommand{\R}{\mathbb{R}}

\title{\textbf{<Course code> <Course title>}\\[4pt]
\large Lecture <N>: <Lecture title>}
\author{<Instructor>}
\date{}

\begin{document}
\maketitle

\section{<Motivation>}
% What students know, and the question this lecture answers. One or two paragraphs.

\section{<Core section>}
% Definitions, ideas, results with short proofs, worked examples, in strict order.
% Add as many core sections as the topic needs, and no more.

% \section{<The method>}   % only if the topic is a method: complete algorithm (+ cost if relevant)

\section{Summary}
\begin{itemize}
\item
\end{itemize}

\section*{Homework <N>}   % only if the course uses homework
\begin{enumerate}
\item
\end{enumerate}

\end{document}
